\documentclass[11pt]{article}
\usepackage[margin=1in]{geometry}
\usepackage{amsmath,amssymb}
\usepackage[hidelinks]{hyperref}

\title{Regularity of the edge ideal of a cycle\\[2pt]
  \large Proof of Conjecture 00000000501}
\author{}
\date{}

\begin{document}
\maketitle

\noindent
Let $n\ge 3$ and $R=k[x_1,\ldots,x_n]$ over an arbitrary field $k$. For
the cycle $C_n$ on the cyclically ordered vertices $1,\ldots,n$, let
$I=I(C_n)$ be the ideal generated by the quadratic monomials corresponding
to its edges. We prove
\[
  \operatorname{reg}_R I=\left\lceil\frac{n-1}{3}\right\rceil+1.
\]

\section*{From Betti numbers to induced subgraphs}

Write $\operatorname{Ind}(G)$ for the simplicial complex of independent
vertex sets of a graph $G$. The quotient $R/I$ is the face ring of
$\operatorname{Ind}(C_n)$. Hochster's formula~\cite[\S5]{hochster} gives
\[
 \beta_{i,j}(R/I)=
   \sum_{\substack{W\subseteq V(C_n)\\ |W|=j}}
   \dim_k \widetilde H_{j-i-1}(\operatorname{Ind}(C_n[W]);k).
\]
Here $C_n[W]$ is the subgraph induced by $W$. Thus a nonzero reduced
homology group in degree $d-1$ of one of these complexes gives a
regularity contribution of $d$ to $R/I$. For $W=\varnothing$, the
complex consists of the empty face and the contribution is zero.

\section*{The full cycle and the proper subsets}

Kozlov's cycle computation~\cite[Proposition~5.2]{kozlov} gives the
homotopy types of $\operatorname{Ind}(C_n)$ as follows.
\[
\begin{array}{c|c|c}
n & \text{homotopy type} & \text{regularity contribution} \\ \hline
3a   & S^{a-1}\vee S^{a-1} & a \\
3a+1 & S^{a-1} & a \\
3a+2 & S^a & a+1
\end{array}
\]
All displayed homology groups in the indicated degree are nonzero over
every field. Therefore the full vertex set contributes
$\lfloor(n+1)/3\rfloor$.

For a nonempty proper subset $W$, the induced subgraph consists of path
components on $l_1,\ldots,l_s$ vertices. Each component is a maximal run
of vertices of $W$ in cyclic order, followed by a distinct vertex outside
$W$. Thus $\sum_{i=1}^s(l_i+1)\le n$. The independence complex is the
simplicial join of the independence complexes of the paths. By the path
computation~\cite[Proposition~4.5]{kozlov}, a path on $l$ vertices has a
contractible independence complex if $l\equiv1\pmod3$ and a sphere of
dimension $\lfloor(l+1)/3\rfloor-1$ otherwise. One can also obtain the path
formula by deleting the third vertex of a path with at least three
vertices. Its neighborhood contains the neighborhood of the first vertex,
so the simplicial map sending vertex three to vertex one and fixing all
other vertices gives a homotopy equivalence to the deletion. The deletion
is a disjoint union of a single edge and the path on $l-3$ vertices, so
its independence complex is the suspension of the latter path complex.

A contractible factor makes a join contractible. If every factor is a
sphere, the join has dimension one less than the sum of their
dimensions shifted by one. Thus the contribution from a proper subset,
when it has nonzero reduced homology, is at most
\[
  \sum_{i=1}^s\left\lfloor\frac{l_i+1}{3}\right\rfloor
  \le \left\lfloor\frac{\sum_{i=1}^s(l_i+1)}{3}\right\rfloor
  \le \left\lfloor\frac{n}{3}\right\rfloor.
\]
The first inequality follows by writing each $l_i+1$ as three times its
quotient plus a nonnegative remainder. It follows that
$\operatorname{reg}_R(R/I)=\lfloor(n+1)/3\rfloor$.

\section*{Regularity of the ideal and Lean verification}

Truncating the minimal free resolution of $R/I$ gives the minimal free
resolution of $I$, so $\beta_{i,j}(I)=\beta_{i+1,j}(R/I)$. The full-cycle
contribution above occurs in homological degree
$n-\lfloor(n+1)/3\rfloor>0$. The shift therefore adds one to the maximum,
giving $\operatorname{reg}_R I=\lfloor(n+1)/3\rfloor+1$. For integer $n$,
$\lfloor(n+1)/3\rfloor=\lceil(n-1)/3\rceil$. This also agrees with the
$s=1$ statement of Beyarslan, H\`a and Trung~\cite[Theorem~5.2]{bht}.

The Lean 4 project states the bidegree nonvanishing criterion from
Hochster's formula using the homology support of the induced complexes.
It proves the run inequality for arbitrary lists, bounds all Betti
shifts, and constructs the extremal bidegree from the full vertex set.
The theorem \texttt{C501.ideal\_regularity} then proves the shift to the
ideal and the ceiling formula from the stated classical inputs. Run
\texttt{lake build} from the \texttt{lean} directory to check it.

\begin{thebibliography}{3}
\bibitem{hochster}
M.~Hochster, \emph{Cohen--Macaulay rings, combinatorics, and simplicial
complexes}, Ring Theory II, Lecture Notes in Pure and Applied Mathematics
\textbf{26} (1977), 171--223.
\href{https://sites.lsa.umich.edu/hochster/wp-content/uploads/sites/1337/2026/03/Cohen-Macaulay-Rings-Combinatorics-and-Simplicial-Complexes.pdf}{Author's copy}.

\bibitem{kozlov}
D.~N. Kozlov, \emph{Complexes of directed trees}, Journal of Combinatorial
Theory, Series A \textbf{88} (1999), 112--122.
\href{https://doi.org/10.1006/jcta.1999.2984}{DOI 10.1006/jcta.1999.2984}.

\bibitem{bht}
S.~Beyarslan, H.~T.~H\`a and T.~N.~Trung, \emph{Regularity of powers of
forests and cycles}, Journal of Algebraic Combinatorics \textbf{42}
(2015), 1077--1095.
\href{https://doi.org/10.1007/s10801-015-0617-y}{DOI 10.1007/s10801-015-0617-y}.
\end{thebibliography}

\end{document}
